\section*{Erd\H{o}s Problem 120: avoiding infinite affine patterns}
\subsection*{Statement and scope}
For every infinite $A\subset\mathbb R$, is there a measurable set $E$ of positive Lebesgue measure containing no set $aA+b$ with $a\ne0$? This is the Erd\H{o}s similarity problem. We prove two classes of infinite patterns can be avoided, explain why finite patterns cannot, and reduce the difficult case to sequences converging to zero. The general question is unresolved here.

The bounded and nowhere-dense constructions retain the valid parts of the earlier GPT-5.2 attempt; the finite-pattern proof and reduction are included to identify the gap precisely. Source statement: \url{https://www.erdosproblems.com/120}. No new resolution is claimed.

\subsection*{Two complete avoidance cases}
If $A$ is unbounded, every $aA+b$ with $a\ne0$ is unbounded, so $E=[0,1]$ avoids it.

Suppose the closure of $A$ contains a nondegenerate interval. Enumerate the rationals of $[0,1]$ as $q_1,q_2,\ldots$ and remove the open intervals of radii $2^{-j-4}$ about them. The remaining set
\[
 E=[0,1]\setminus\bigcup_{j\ge1}(q_j-2^{-j-4},q_j+2^{-j-4})
\]
is closed, has empty interior, and has measure at least $1-\sum_{j\ge1}2^{-j-3}=7/8$. If $aA+b\subset E$, closedness gives $a\overline A+b\subset E$, contradicting its empty interior. Thus this $E$ avoids every affine copy of $A$.

\subsection*{Why the infinite hypothesis matters}
Every measurable set $E$ of positive measure contains an affine copy of every finite set $F=\{u_1,\ldots,u_k\}$. Choose a bounded measurable $E_0\subset E$ of finite positive measure. Translations are continuous in $L^1$:
\[
 \mu((E_0-h)\mathbin{\triangle}E_0)\longrightarrow0\quad(h\to0).
\]
One proof approximates $E_0$ in measure by a finite union of bounded intervals, for which the assertion follows by adding the changes at finitely many endpoints.

Choose a sufficiently small nonzero $a$ so that
$\mu(E_0\setminus(E_0-au_i))<\mu(E_0)/(2k)$ for every $i$. The union bound gives
\[
 \mu\left(E_0\cap\bigcap_{i=1}^k(E_0-au_i)\right)\ge\mu(E_0)/2>0.
\]
Any $b$ in this intersection satisfies $b+aF\subset E$. The same estimate cannot be applied to an infinite collection with a fixed nonzero $a$.

\subsection*{A reduction of the remaining case}
Every infinite bounded subset of $\mathbb R$ has a sequence of distinct points converging to a limit $c$. Passing to points on one side of $c$ and then to a subsequence, their distances from $c$ can be made strictly decreasing to zero. Translation and, if necessary, reflection turn that subsequence into
\[
 b_1>b_2>\cdots>0,\qquad b_j\to0.
\]
If some positive-measure set avoids every affine copy of this normalized sequence, it avoids every affine copy of the original $A$: any such copy of $A$ would contain the corresponding copy of the subsequence. Together with the unbounded case, this shows that it suffices to settle all strictly decreasing positive sequences tending to zero.

\subsection*{Remaining gap}
A convergent sequence can lie in a closed nowhere-dense set, so the second construction supplies no contradiction for this reduced case. The finite-pattern argument gives potentially different, shrinking scales for different prefixes; it supplies no common nonzero scale for an infinite pattern. Even the geometric-sequence case is not established here.
\subsection*{COMPLETION ESTIMATE}
\textbf{COMPLETION: 30\%}
